\documentclass[fontset=ubuntu]{article}

% Chinese support
\usepackage{ctex}

% customize section/subsection title
\usepackage{titlesec}

% customize sep of items
\usepackage{enumitem}

\usepackage{fontawesome}

% define some colors
\usepackage[x11names]{xcolor}

\usepackage{hyperref}

\hypersetup{colorlinks=true, urlcolor=Red4}

\setitemize[1]{ itemsep=0.3ex, % space between item
partopsep=0.0ex, %
parsep=\parskip, topsep=0.3ex % space before text
}

% set margin
\usepackage[
 a4paper,
 left=0.8in,
 right=0.8in,
 top=0.70in,
 bottom=0.55in,
 nohead
]{geometry}

% do not need indent in resume
\setlength{\parindent}{0pt}

% ----------------------------------------------------------------------------- %

\titleformat{\section}
% format: to be applied to the whole title -- label and text
{\color{Red4}\Large\bf\raggedright}
% label: prefix the title body
{}
% sep: space between label and title body
{0em}
% before: code preceding the title body
{}
% after: code following the title body
% insert a horizontal line after the heading
[\titlerule]

% adjust the vertical space before and after the heading
\titlespacing*{\section}{0cm}{*1.6}{*1.2}

% similar above
\titleformat{\subsection}
{\normalsize\raggedright}
{}{0em}{}
\titlespacing*{\subsection}{0cm}{*1.6}{0.0ex}

% ----------------------------------------------------------------------------- %

\newcommand{\dateditem}[2]{% set the date text Red4
\subsection[#1]{#1 \hfill \textcolor{Red4}{#2}}
}

\makeatletter
\renewcommand\maketitle{
 \begin{center}
  \LARGE\textbf{\@title}
 \end{center}
}
\makeatother

\title{吴坎}

% ----------------------------------------------------------------------------- %

\begin{document}
 \maketitle

 \begin{center}
  \texttt{
   \faLink ~ \url{https://wu-kan.cn} 
   \textperiodcentered\
   \faGithub ~ github.com/wu-kan 
   \textperiodcentered\
   \faEnvelope ~ cv@wu-kan.cn 
   \textperiodcentered\
   \faWechat ~ wukan0621
  }
 \end{center}

 % ----------------------------------------------------------------------------- %

 \section{Education}

 \dateditem{\textbf{中山大学} \quad 计算机科学与技术 \quad 本科 \textperiodcentered\ 硕博连读}{2017年9月 -- 2026年6月（预期）}

 \begin{itemize}
  \item 导师：\href{https://cse.sysu.edu.cn/content/2483}{卢宇彤}、\href{https://xianweiz.github.io/}{张献伟}@\href{http://nscc-gz.cn/}{国家超级计算广州中心}
  \item 研究方向：异构系统设计与加速；软硬件结合的体系优化
 \end{itemize}

 \section{Project}

 \dateditem{\textbf{\large SMILE: LLC-based Shared Memory Expansion to Improve GPU Thread Level Parallelism}}{2022年4月 -- 2023年6月}
 
 Tianyu Guo, Xuanteng Huang, \textbf{Kan Wu}, Xianwei Zhang, Nong Xiao. \href{https://www.dac.com/Conference/2024-Speaker-Resource-Page}{DAC'2024}
 
 \begin{itemize}
  \item 提出面向 GPU 架构的硬件创新 SMILE，使用 Last Level Cache 提高线程并行度

  \item 在项目前期提出初版硬件架构设计
 \end{itemize}

 \dateditem{\textbf{\large RollBin: Reducing code-size via loop rerolling at binary level}}{2022年2月 -- 3月}
 Tianao Ge, Zewei Mo, \textbf{Kan Wu}, Xianwei Zhang, Yutong Lu.
 \href{https://pldi22.sigplan.org/details/LCTES-2022/9/RollBin-Reducing-code-size-via-loop-rerolling-at-binary-level}{LCTES'22}
 \begin{itemize}
  \item 提出一种二进制代码优化手段 RollBin，在实验中相比已经存在的技术提升 31\%
    -- 38\%
  \item 参与论文思路讨论、实验验证、数据整理与可视化、论文撰写
 \end{itemize}

 \dateditem{\textbf{\large 面向新一代神威超级计算机的Vision Transformer模型训练过程优化}}{2021年8月 -- 9月}

 SACA（Sunway Accelerate Computing Architecture.
 \href{https://sysu-scc.feishu.cn/file/boxcnsEncoYfL79hceRLrAP11Jg}{“神威杯”参赛作品}

 \begin{itemize}
  \item 在新一代神威超算上优化 SWPyTorch 使用 256 节点训练 ViT 模型的过程
  \item 前向过程加速上千倍，反向与权重更新过程加速上百倍
  \item 算例规模下六核组 SGEMM 峰值性能超过 8.2 TFlops，达到理论性能的 59.24\%
 \end{itemize}

 \dateditem{\textbf{\large 面向国产异构处理器的HPL-AI基准实现及优化}}{2020年11月 -- 2021年2月}

 AscendCL/CUDA. \href{https://wu-kan.cn/2021/03/14/HPL-AI/}{本科毕业设计}

 \begin{itemize}
  \item 首个开源的多节点 HPL-AI 基准实现，基于 hpl-2.3 并达到 2.59x 加速比

  \item 使用 NPU（昇腾910）的半精度算力加速 Linpack 求解，然后通过数值迭代提升精度

  \item 在X86、ARM、GPU、NPU等多种环境上通过正确性、可扩展性验证
 \end{itemize}

 % ----------------------------------------------------------------------------- %

 \section{Experience}

 \dateditem{\textbf{中山大学}\quad 计算机学院}{2021年9月 --2022年6月}
 
 首届\href{https://cse.sysu.edu.cn/content/5366}{``冯''班}
 助教. 编译原理, 组成原理
 
 \begin{itemize}
  \item 独立设计编译原理实验 \href{https://github.com/arcsysu/SYsU-lang}{SYsU-lang}，并应用于课堂教学

  \item 发起并参与设计组成原理实验 \href{https://yatcpu.sysu.tech/}{YatCPU}（“逸芯”）
 \end{itemize}

 \dateditem{\textbf{Bytedance}\quad AI Lab Speech \& Audio Team, Beijing}{2021年2 -- 3月，2021年7 -- 8月}
 
 算法引擎实习生. Panther 推理引擎 GPU 后端
 
 \begin{itemize}
  \item 开发定制 TransformerDecoder 算子，同时编写对应单元测试、图匹配脚本
   \begin{itemize}
    \item 相关模型在 NVIDIA T4 上推理加速$1.6\times$（80ms$\to$50ms）
   \end{itemize}

  \item 基于CUTLASS的量化multihead attention算子（batchGemm-Softmax-batchGemm）
   \begin{itemize}
    \item 序列长度较短时，相较基于cublasLt的原始实现加速7倍以上
   \end{itemize}
 \end{itemize}

 \dateditem{\textbf{中山大学}\quad 国家超级计算广州中心}{2020年3月 -- 2021年6月}

 队长. \href{https://scc.sysu.tech}{中山大学超算队}

 \begin{itemize}
  \item 在 ASC 决赛中
   \begin{itemize}
    \item 调整主板接线实现激进功耗控制，\href{https://mp.weixin.qq.com/s?__biz=MzI5NzUxMDAxOQ==&mid=2247486028&idx=1&sn=08b70942a6e7b0f8e4cd42a98789f4fe}{在三千瓦功耗限制下打破
     Linpack 纪录并提升57\%}

    \item \href{https://wu-kan.cn/2021/05/11/\%E9\%87\%8F\%E5\%AD\%90\%E7\%BA\%BF\%E8\%B7\%AF\%E6\%A8\%A1\%E6\%8B\%9F\%E5\%99\%A8QuEST\%E5\%9C\%A8\%E5\%A4\%9AGPU\%E5\%B9\%B3\%E5\%8F\%B0\%E4\%B8\%8A\%E7\%9A\%84\%E6\%80\%A7\%E8\%83\%BD\%E4\%BC\%98\%E5\%8C\%96/}{负责量子线路模拟器赛题QuEST在多GPU平台上的性能优化}
   \end{itemize}

  \item 在队内定期开展\href{https://scc.sysu.tech/Talks/}{技术分享}，在国内所有超算相关竞赛获前三名，受邀
   IPCC 赛前讲座
   \begin{itemize}
    \item \href{https://sysu-scc.feishu.cn/file/boxcnyn3xp9JJz7rNyjUtUAhHLc}{超算竞赛备赛指南}，\href{https://sysu-scc.feishu.cn/file/boxcnjtWoYlkvSF4pF4grCwudFr}{超算基本原理与方法论}，\href{https://sysu-scc.feishu.cn/file/boxcnQcRimOByBFREJEc5lgtgXc}{现代处理器优化概论}
   \end{itemize}
 \end{itemize}

 % ----------------------------------------------------------------------------- %

 \section{Honor}
 \dateditem{\href{https://mp.weixin.qq.com/s/ahnG7gJLBu9CMCYjk1_CkQ}{\texttt{二等奖\& oneAPI技术创新一等奖}}，“英特尔杯”并行应用挑战赛（PAC），应用组}{2022 年 12 月}
 \dateditem{\href{https://mp.weixin.qq.com/s/k7mjyaKs-rGPOcMQTIrKqg}{\texttt{金奖（1/unknown）}}，鲲鹏应用创新大赛，HPC高性能计算赛道}{2022 年 10 月}
 \dateditem{\texttt{\href{https://github.com/arcsysu/SYsU-lang-paper}{优秀论文一等奖}}，中国高校计算机教育大会（CCEC）}{2022 年 7 月}
 \dateditem{\href{https://mp.weixin.qq.com/s/L7YX4aaaNAlNCKDYOEhMSA}{\texttt{二等奖（3/130+）}}，ACM中国-国际并行挑战赛（ACM-IPCC）}{2021 年 11 月}
 \dateditem{\href{https://mp.weixin.qq.com/s/GQsaR_V1NoGzCvbGARMeEA}{\texttt{二等奖（3/100+）}}，“神威杯”并行应用挑战赛（CPC）}{2021 年 10 月}
 \dateditem{\href{https://mp.weixin.qq.com/s/AJKqG2KvLl6n_88_8IzsYw}{\texttt{CCF优秀大学生（全国共94名，中山大学共1名）}}，中国计算机学会（CCF-ECA）}{2021 年 9 月}
 \dateditem{\href{https://wu-kan.cn/2021/03/14/HPL-AI/}{\texttt{优秀本科毕业论文（Top 5\%）}}，中山大学}{2021 年 5 月}
 \dateditem{\href{https://wu-kan.cn/2021/05/19/\%E6\%88\%91\%E7\%9A\%84ASC\%E5\%86\%B3\%E8\%B5\%9B\%E5\%A4\%8D\%E7\%9B\%98-\%E5\%86\%92\%E9\%99\%A9-\%E5\%A4\%B1\%E8\%AF\%AF\%E4\%B8\%8E\%E7\%BF\%BB\%E8\%BD\%A6/}{\texttt{最高性能奖 \& 一等奖（3/400+）}}，世界大学生超级计算机竞赛（ASC）}{2021 年 5 月}
 \dateditem{\href{https://wu-kan.cn/_posts/2020-10-17-2020-CCF-CCSP\%E7\%AB\%9E\%E8\%B5\%9B-\%E5\%90\%AB\%E5\%88\%86\%E8\%B5\%9B\%E5\%8C\%BA\%E7\%AB\%9E\%E8\%B5\%9B/}{\texttt{金奖（4/98）}}，CCF 大学生计算机系统与程序设计竞赛（CCSP），华南赛区}{2020 年 10 月}
 \dateditem{\texttt{一等奖学金\& 学科竞赛一等奖学金（Top 5\%）}，中山大学}{2020 年 10 月}
 \dateditem{\href{https://mp.weixin.qq.com/s/ahnG7gJLBu9CMCYjk1_CkQ}{\texttt{银牌（2/100+）}}，“英特尔杯”并行应用挑战赛（PAC），优化组}{2020 年 9 月}
 \dateditem{\href{http://cse.sysu.edu.cn/content/5412}{\texttt{一等奖（3/unknown）}}， 中科院“先导杯”并行应用挑战赛（CAS-PRA）}{2020 年 8 月}
 \dateditem{\href{https://wu-kan.cn/_posts/2019-12-16-\%E7\%AC\%AC\%E5\%8D\%81\%E5\%85\%AB\%E6\%AC\%A1CCF\%E8\%AE\%A1\%E7\%AE\%97\%E6\%9C\%BA\%E8\%BD\%AF\%E4\%BB\%B6\%E8\%83\%BD\%E5\%8A\%9B\%E8\%AE\%A4\%E8\%AF\%81/}{\texttt{Top 0.05\%（6/11395）}}，CCF 计算机软件能力认证（CCF-CSP）}{2019 年 12 月}
 \dateditem{\href{https://wu-kan.cn/_posts/2019-11-04-\%E5\%86\%8D\%E8\%A7\%81-\%E7\%AE\%97\%E6\%B3\%95\%E7\%AB\%9E\%E8\%B5\%9B/}{\texttt{银牌}}，国际大学生程序设计竞赛（ICPC），徐州站}{2019 年 11 月}
 \dateditem{\texttt{金奖$\times 2$}，中国大学生程序设计竞赛（CCPC），广东赛区}{{\href{https://ccpc.io/post/110}{2018 年 5 月}，\href{http://www.gdcpc.cn/ArticleView/ArticleDetail?articleId=5cda35a3659ce607a808929a}{2019 年 5 月}}}

\end{document}